\lesson{4}{Dec 01 2021 Wed (18:08:39)}{Solving and Graphing PF}{Unit 3}

A function written in factored form gives us the roots of the equation. The roots and the end behavior rules below can be used to create a graph for the function.

Some \bf{Polynomial Functions} are far more complicated and should be graphed using graphing technology such as graphing software or a graphing calculator. The real number solutions of a \bf{Polynomial Equation} can be found by identifying the \bf{x-intercepts} from the graph of the function. Graph the function, using graphing technology and find the point(s) of intersection between the graph and the x-axis.

\begin{marginfigure}
    \centering
    \incfig{graph-technology}
    \sidecaption{$0.8x^{4} \times 8x^{3} + 2x^{2} - 7x - 1$ Graphed.}
    \label{fig:graph-technology}
\end{marginfigure}

\begin{definition}[Average Rate of Change]
    \begin{align}
        m = \frac{y_2 - y_1}{x_2 - x_1}
    \end{align}
    
    \paragraph{Increasing versus Decreasing}
    
    \begin{itemize}
        \item
        \item If the \bf{Average Rate of Change} is \bf{Positive}, then the function is considered to be \bf{Increasing}.
        \item If the \bf{Average Rate of Change} is \bf{Negative}, then the function is considered to be \bf{Decreasing}.
    \end{itemize}
    
    \paragraph{Minimums and Maximums}
    
    \begin{itemize}
        \item
        \item \bf{Minimums} are located on the \bf{Lowest Point of a Graph}. A \bf{Local Minimum} is the bottom of a \bf{Valley} or \bf{Turning Point} where the graph goes from \bf{Decreasing} to \bf{Increasing}.
        \item \bf{Maximums} are located on the \bf{Highest Point of a Graph}. A \bf{Local Maximum} is the top of a \bf{Hill} or \bf{Turning Point} where the graph goes from \bf{Increasing} to \bf{Decreasing}.
    \end{itemize}
\end{definition}

\begin{example}[Example 1]
    Determine the zeros of the function $f(x) = (x - 2)(x - 5)(x + 1)$, and describe the end behavior of the graph.
    The zeros of the function represent the values of $x$ that make the function equal to $0$. Replace the function notation with $0$ to determine the zeros of the function.
    
    \begin{align}
        f(x) &= (x - 2)(x - 5)(x + 1) \\
        0 &= (x - 2)(x - 5)(x + 1) \\
        0 &= x^3 - 6 \times x^2 + 3x + 10 \\
        0 &= x^3 - 6 \times x^2 + 3x + 10 \\
        x &= 0, x = 5, x = -1 \\
    \end{align}
    
    So, from this, we know that the zeros function from this function are: $0, 5, -1$, which means that's where they will intersect the $x-axis$.
    
    The function has an \bf{Odd Degree}, so the ends will travel in the opposite directions. The leading coefficient is positive, so the left side continues down and the right side continues up.
\end{example}

\newpage
